\documentclass[12pt]{article}
\usepackage[utf8]{inputenc}
\usepackage[T1]{fontenc}
\usepackage{hyperref}
\title{Kasem Language Resources for K-12 Curriculum Development}
\author{Compiled by Claude Sonnet 5.5 (claude-sonnet-5-5) for Andr\'as Kornai}
\date{}
\begin{document}
\maketitle

\section{Introduction}
Kasem (Kas\'im, Kasena; ISO 639-3 xsm) is a Gur language spoken in Burkina Faso and Ghana. Speaker estimates vary, from about 250,000 in total to about 120,000 in each country. Most published material comes from SIL and the Ghana Institute of Linguistics, Literacy and Bible Translation (GILLBT).

\section{Primary Educational Resources}
Search results listed these GILLBT items (bibliographic mentions, no direct URLs seen):
\begin{itemize}
\item \emph{De lage de zamse kasem 2}, a primer, GILLBT, Tamale, 1984.
\item \emph{Let's Write Kasem}, a spelling guide, GILLBT, 2002.
\end{itemize}

\section{Linguistic and Reference Materials}
\begin{itemize}
\item Kasem--French--English dictionary (Burkina Faso) by Urs and Idda Niggli, online at Webonary: \url{https://www.webonary.work/kassem-bf/?p=9444}
\item Dictionary PDF, volume 1 (A--K): \url{https://www.webonary.work/kassem-bf/files/Dictionnaire-Kassem-francais-anglais-A-K.pdf}
\item Dictionary foreword: \url{https://www.webonary.work/kassem-bf/overview/forward/?lang=en}
\item Kasem dictionary (Ghana), based on lexical files by John and Kathleen Callow: \url{https://www.webonary.work/kasem-gh/?p=146819}
\item Language page of the Ghana dictionary: \url{https://www.webonary.work/kasem-gh/language/}
\item Philip Hewer, \emph{A Basic Grammar of Kasem} (1981), Glottolog reference: \url{https://glottolog.org/resource/reference/id/29269}
\end{itemize}

\section{Religious and Community Resources}
A Kasem Bible translation committee in Ti\'eb\'el\'e and P\^o is credited in the dictionary materials. A French-language page reports a Kasem New Testament for smartphones with read-and-listen format, but no URL for it was seen. Joshua Project profile (Ghana): \url{https://www.joshuaproject.net/index.php/people_groups/12556/GH}

\section{Digital and Audio Resources}
The Webonary dictionaries are online. The smartphone New Testament mentioned above is the only audio-related resource reported.

\section{Curriculum Development Context}
Kasem is used in some schools, but no curriculum documents were found in search results. A curriculum could build on the GILLBT primer and spelling guide, the two bilingual dictionaries and the Hewer grammar. Orthographies and dialects differ between Burkina Faso and Ghana, which would need reconciling or separate tracks. Contact GILLBT and SIL Burkina Faso for current classroom materials.

\end{document}